\clearpage
\renewcommand{\sectioncolor}{alert}
\section{Where Monte Carlo breaks --- measured}
\markboth{Where Monte Carlo breaks}{}

Now the method everyone actually uses on systems that are \emph{not} exactly solvable: Metropolis.
Propose a spin flip, accept with probability $\min(1,e^{-\beta\Delta E})$, run long enough.

\textbf{That last phrase is where everything goes wrong}, and the Ising model lets us measure exactly
how wrong.

\begin{keyidea}[title={A perfect test, because we know the answer exactly}]
At $h=0$ the model is exactly symmetric under flipping every spin. Therefore
\[
\langle M\rangle = 0 \qquad\text{exactly, for every }\beta,\text{ every }N.
\]
No approximation, no limit. So: run several \emph{independent} chains and see whether they all return
zero. \textbf{Every chain below runs the identical, correct algorithm. None of them is buggy.}
\end{keyidea}

\begin{center}
\begin{tikzpicture}
\begin{axis}[width=15.2cm,height=6.6cm,
  xlabel={\footnotesize $\beta=1/T$},
  ylabel={\footnotesize per-chain $\langle M\rangle/N$},
  xmin=0.3,xmax=3.25,ymin=-1.25,ymax=1.25,
  xtick={0.5,1.0,1.5,2.0,2.5,3.0}, ytick={-1,-0.5,0,0.5,1},
  tick label style={font=\scriptsize}, label style={font=\footnotesize},
  grid=major, grid style={slate!18,line width=0.3pt}, axis line style={slate!60}, clip=false]
 \addplot[only marks,mark=*,mark size=2.1pt,color=alert] coordinates {
   (0.5,-0.002)(0.5,0.010)(0.5,0.008)(0.5,0.009)(0.5,-0.001)(0.5,-0.011)
   (1.0,-0.021)(1.0,-0.013)(1.0,0.029)(1.0,0.014)(1.0,-0.037)(1.0,0.010)
   (1.5,0.121)(1.5,-0.021)(1.5,-0.078)(1.5,0.126)(1.5,-0.026)(1.5,0.049)
   (2.0,0.116)(2.0,0.486)(2.0,-0.749)(2.0,0.707)(2.0,0.311)(2.0,0.962)
   (2.5,-0.162)(2.5,-0.950)(2.5,0.953)(2.5,-0.448)(2.5,-0.982)(2.5,-0.999)
   (3.0,1.000)(3.0,-1.000)(3.0,-0.993)(3.0,-1.000)(3.0,-1.000)(3.0,-1.000)};
 \draw[bio,line width=1.6pt] (axis cs:0.3,0)--(axis cs:3.25,0);
 \node[font=\scriptsize\bfseries,text=bio,anchor=south west] at (axis cs:0.35,0.04)
   {the true answer, exactly: $\langle M\rangle=0$};
 \draw[slate!50,dashed,line width=0.8pt] (axis cs:1.75,-1.25)--(axis cs:1.75,1.25);
 \node[font=\scriptsize\bfseries,text=bio,anchor=north] at (axis cs:1.0,-1.05) {chains agree};
 \node[font=\scriptsize\bfseries,text=alert,anchor=north] at (axis cs:2.55,-1.05)
   {chains disagree with each other \emph{and} with the truth};
 \node[font=\scriptsize,text=slate,anchor=south] at (axis cs:2.0,1.08) {$\xi=27$};
 \node[font=\scriptsize,text=slate,anchor=south] at (axis cs:2.5,1.08) {$\xi=74$};
 \node[font=\scriptsize,text=slate,anchor=south] at (axis cs:3.0,1.08) {$\xi=202$};
\end{axis}
\end{tikzpicture}
\end{center}
\vspace{2pt}
\begin{center}\begin{minipage}{0.93\textwidth}\footnotesize\color{slate}
\textbf{Figure 5.}\; Six independent Metropolis chains, $N=64$, 6{,}000 sweeps each. Measured for
this dossier. At $\beta=3$ the chains report $\pm1$ when the exact answer is $0$.
\end{minipage}\end{center}

\begin{center}\small
\begin{tabular}{@{}rrrl@{}}
\toprule
\rowcolor{alert!12}
$\beta$ & $\xi$ & \textbf{spread across chains} & \textbf{verdict}\\
\midrule
0.5 & 1.30 & 0.021 & \textcolor{bio}{\textbf{ok}}\\
1.0 & 3.67 & 0.066 & \textcolor{bio}{\textbf{ok}}\\
1.5 & 10.03 & 0.203 & \textcolor{sand!80!black}{\textbf{poor}}\\
2.0 & 27.30 & 1.710 & \textcolor{alert}{\textbf{BROKEN}}\\
2.5 & 74.21 & 1.952 & \textcolor{alert}{\textbf{BROKEN}}\\
\rowcolor{alert!8}
3.0 & 201.71 & \textbf{2.000} & \textcolor{alert}{\textbf{BROKEN}}\\
\bottomrule
\end{tabular}
\end{center}
\vspace{-2pt}
\begin{center}\begin{minipage}{0.93\textwidth}\footnotesize\color{slate}
\textbf{Table 3.}\; A spread of 2.000 means some chains returned $+1$ and others $-1$ --- the maximum
possible disagreement. An independent 8-chain run at 30{,}000 sweeps reproduced the same pattern.
\end{minipage}\end{center}

\begin{missing}[title={Why --- and it is the same mechanism as the whole of Part IX}]
As $\beta$ grows, $\xi$ grows exponentially (Table 2). Once $\xi\gtrsim N$ the entire ring is a
single domain, and to get from ``mostly up'' to ``mostly down'' the chain must pass through a
high-energy configuration. The waiting time for that is $\sim e^{\beta\Delta E}$ ---
\textbf{exponential in the barrier height}.
\begin{center}
\begin{tikzpicture}[font=\scriptsize,
  n/.style={rounded corners=3pt,draw=#1,line width=0.9pt,fill=#1!8,inner sep=6pt,
            text width=3.3cm,align=center,minimum height=1.25cm}]
 \node[n=nano]  (a) at (0,0)     {temperature falls};
 \node[n=sand]  (b) at (4.0,0)   {spectral gap closes\\$\lambda_+/\lambda_-\to1$};
 \node[n=ember] (c) at (8.0,0)   {$\xi$ grows exponentially};
 \node[n=alert] (d) at (12.3,0)  {barrier crossing time $\sim e^{\beta\Delta E}$\\\textbf{chain never mixes}};
 \foreach \x/\y in {a/b,b/c,c/d}{\draw[-{Stealth[length=4.5pt]},line width=1pt,draw=slate] (\x)--(\y);}
\end{tikzpicture}
\end{center}
\vspace{-2pt}
This is the quantitative form of the sentence in Part IX §2: \emph{``Monte Carlo converges at
$O(n^{-1/2})$ provided it mixes. On a rugged landscape with barriers $\gg k_BT$, it does not mix.''}
And it is the argument for Boltzmann generators: rather than \emph{crossing} the barrier, learn a
change of variables in which there is no barrier to cross.
\end{missing}

\begin{plain}
The lesson is not ``Monte Carlo is bad''. It is that a method can be provably correct, correctly
implemented, and still give you a confidently wrong answer --- because ``run it long enough'' hid an
exponential. You would never know, in a system where the true answer was not available. That is the
situation in every protein simulation ever run.
\end{plain}
